
\documentclass[12pt,letterpaper]{article}

\usepackage[spanish,es-tabla]{babel}
\usepackage[utf8]{inputenc}
\usepackage[T1]{fontenc}
\usepackage{amsmath, amssymb, amsfonts}
\usepackage{geometry}
\usepackage{graphicx}
\usepackage{xcolor}
\usepackage{enumitem}
\usepackage{multicol}
\usepackage{fancyhdr}
\usepackage[most]{tcolorbox}
\usepackage{hyperref}
\usepackage{transparent}
\usepackage{eso-pic}
\usepackage{array}

\geometry{left=2cm,right=2cm,top=3cm,bottom=2.35cm,headheight=40pt,headsep=0.65cm,footskip=1cm}

% =========================================================
% INFORMACIÓN DE LA GUÍA
% =========================================================
\newcommand{\institucion}{Usach Premium}
\newcommand{\curso}{Álgebra II}
\newcommand{\guiasnumero}{Guía 3 - Espacios Vectoriales y Transformaciones Lineales}
\newcommand{\semestre}{1er. Semestre de 2026}
\newcommand{\preparadopor}{Francisca Varela}
\newcommand{\webinstitucion}{www.usachpremium.cl}
\newcommand{\instagraminstitucion}{@usach.premium}
\newcommand{\instagramurl}{https://www.instagram.com/usach.premium}
\newcommand{\logoup}{logo_up.png}
\newcommand{\logousach}{logo_usach.png}
\newcommand{\logofondo}{logo_up.png}
\newcommand{\transparencia}{0.055}
\newcommand{\fraseportada}{La práctica constante transforma el procedimiento en claridad.}
\newcommand{\listacontenidos}{
    \item[\color{azulFuerte}$\Rightarrow$] Espacios vectoriales y subespacios.
    \item[\color{azulFuerte}$\Rightarrow$] Sistemas de generadores.
    \item[\color{azulFuerte}$\Rightarrow$] Independencia lineal, bases y dimensión.
    \item[\color{azulFuerte}$\Rightarrow$] Transformaciones lineales, núcleo e imagen.
}

% =========================================================
% COLORES
% =========================================================
\definecolor{celesteSuave}{HTML}{F0F7FF}
\definecolor{azulFuerte}{HTML}{1A5276}
\definecolor{turquesaUsach}{HTML}{33CCCC}
\definecolor{enlace}{HTML}{1A5276}

\hypersetup{
    colorlinks=true,
    linkcolor=enlace,
    urlcolor=enlace
}

% =========================================================
% ENCABEZADO Y PIE
% =========================================================
\pagestyle{fancy}
\fancyhf{}
\lhead{\includegraphics[height=1.1cm]{\logoup}}
\chead{\small \textbf{\institucion}\\ \curso \\ \href{\instagramurl}{\instagraminstitucion}}
\rhead{\includegraphics[height=1.1cm]{\logousach}}
\lfoot{\small \fraseportada}
\rfoot{\small Página \thepage}
\renewcommand{\headrulewidth}{0.4pt}
\renewcommand{\footrulewidth}{0.4pt}

\newcommand{\MarcaAgua}{
    \AddToShipoutPicture{
        \AtPageCenter{
            \makebox(0,0){%
                \transparent{\transparencia}%
                \includegraphics[width=0.7\textwidth]{\logofondo}%
            }
        }
    }
}

% =========================================================
% CAJAS
% =========================================================
\tcbset{
    colback=celesteSuave,
    colframe=azulFuerte,
    arc=2mm,
    boxrule=0.7pt
}

\newtcolorbox{kitpropiedades}[1]{
    title=\textbf{#1},
    coltitle=white,
    colbacktitle=azulFuerte,
    fonttitle=\bfseries
}

\newtcolorbox{desafiobox}[1]{
    title=\textbf{#1},
    colback=white,
    colframe=azulFuerte,
    coltitle=white,
    colbacktitle=azulFuerte,
    fonttitle=\bfseries,
    drop shadow
}

\newtcolorbox{solbox}[1]{
    title=\textbf{#1},
    coltitle=white,
    colbacktitle=azulFuerte,
    colback=celesteSuave,
    colframe=azulFuerte,
    fonttitle=\bfseries,
    breakable,
    enhanced,
    before skip=7pt,
    after skip=9pt
}

\newtcolorbox{conclusionbox}{
    colback=white,
    colframe=azulFuerte,
    boxrule=0.6pt,
    arc=1.5mm,
    left=2mm, right=2mm, top=1mm, bottom=1mm,
    before skip=5pt,
    after skip=2pt
}

\newenvironment{introduccion}{
    \MarcaAgua
    \section*{Introducción}
}{}

\newenvironment{ejercicios}{
    \MarcaAgua
    \section*{Ejercicios}
}{}

\newenvironment{soluciones}{
    \newpage
    \MarcaAgua
    \begin{tcolorbox}[colback=azulFuerte,colframe=azulFuerte]
        \centering \Large \textcolor{white}{\textbf{SOLUCIONARIO}}
    \end{tcolorbox}
    \small
}{}

% =========================================================
% DOCUMENTO
% =========================================================

\begin{document}
\thispagestyle{plain}

\begin{center}
    \vspace*{0.2cm}
    \begin{minipage}{0.45\textwidth}
        \raggedright
        \includegraphics[height=1.7cm]{\logoup}
    \end{minipage}
    \begin{minipage}{0.45\textwidth}
        \raggedleft
        \includegraphics[height=1.7cm]{\logousach}
    \end{minipage}

    \vspace{3.1cm}

    {\Large \textbf{\institucion}}\\[0.3cm]
    {\Huge \textbf{\curso}}\\[0.35cm]
    {\Large \textbf{\guiasnumero}}\\[0.6cm]
    {\large \textit{Unidad III: Espacios Vectoriales y Transformaciones Lineales}}\\[1.2cm]

    \begin{tcolorbox}[width=0.78\textwidth, colback=celesteSuave, colframe=azulFuerte, title=\textbf{Contenidos de la guía}]
        \begin{itemize}[leftmargin=1.2cm]
            \listacontenidos
        \end{itemize}
    \end{tcolorbox}

    \vfill
    \begin{minipage}{0.45\textwidth}
        \raggedright
        \textbf{\institucion}\\
        \semestre\\
        \textbf{Preparado por: \preparadopor}
    \end{minipage}
    \begin{minipage}{0.45\textwidth}
        \raggedleft
        Web: \href{https://usachpremium.cl/}{\webinstitucion}\\
        Instagram: \href{\instagramurl}{\instagraminstitucion}
    \end{minipage}
\end{center}

\newpage

\begin{introduccion}
Esta guía está diseñada en modo ascenso: comienza con ejercicios de verificación directa de subespacios, continúa con sistemas de generadores, bases y dimensión, y finaliza con problemas de transformaciones lineales tipo control. Se recomienda justificar cada respuesta indicando la definición o teorema utilizado y, cuando corresponda, resolver el sistema lineal asociado.

\begin{kitpropiedades}{Kit de propiedades útiles}
\begin{multicols}{2}
\begin{itemize}[leftmargin=0.7cm]
    \item Subespacio: \(W\leq V\) si contiene al vector cero, es cerrado para la suma y para la ponderación escalar.
    \item Generador: \(W=\langle v_1,\ldots,v_k\rangle\) si todo vector de \(W\) se escribe como combinación lineal de ellos.
    \item Independencia lineal: \(\alpha_1v_1+\cdots+\alpha_kv_k=0\Rightarrow \alpha_1=\cdots=\alpha_k=0\).
    \item Base: conjunto linealmente independiente que genera el espacio.
    \item Dimensión: número de vectores de una base.
    \item Núcleo: \(\ker(T)=\{v\in V:T(v)=0\}\).
    \item Imagen: \(\operatorname{Im}(T)=\{T(v):v\in V\}\).
    \item Teorema núcleo-imagen: \(\dim(\ker T)+\dim(\operatorname{Im}T)=\dim(V)\).
    \item Isomorfismo: una transformación lineal \(T:V\to W\) es isomorfismo si es biyectiva, es decir, inyectiva y sobreyectiva.
\end{itemize}
\end{multicols}
\end{kitpropiedades}
\end{introduccion}

\begin{ejercicios}

\begin{enumerate}[leftmargin=0.65cm,label=\textbf{\arabic*.}]

\item Determine si
\[
W=\{(x,y,z)\in\mathbb{R}^3:x-2y+z=0\}
\]
es un subespacio de \(\mathbb{R}^3\).

\item Determine si
\[
W=\{(x,y,z)\in\mathbb{R}^3:x+y=1\}
\]
es un subespacio de \(\mathbb{R}^3\).

\item Determine si
\[
W=\left\{
\begin{pmatrix}
a & b\\
c & d
\end{pmatrix}\in M_{\mathbb{R}}(2):a+d=0
\right\}
\]
es un subespacio de \(M_{\mathbb{R}}(2)\).

\item Determine si
\[
W=\{a_0+a_1x+a_2x^2\in\mathbb{R}_2[x]:a_0-a_1+a_2=0\}
\]
es un subespacio de \(\mathbb{R}_2[x]\).

\item Sea
\[
W=\{(x,y,z)\in\mathbb{R}^3:x+y-z=0\}.
\]
Determine un sistema generador para \(W\).

\item Sea
\[
W=\{(x,y,z,t)\in\mathbb{R}^4:x-y+t=0\}.
\]
Determine un sistema generador para \(W\).

\item Sea
\[
W=\left\{
\begin{pmatrix}
a & b\\
c & d
\end{pmatrix}\in M_{\mathbb{R}}(2):a=d
\right\}.
\]
Determine un sistema generador para \(W\).

\item Sea
\[
W=\{p(x)\in\mathbb{R}_2[x]:p(1)=0\}.
\]
Determine un sistema generador para \(W\).

\item Determine si el conjunto
\[
\{(1,2),(2,4)\}
\]
es linealmente independiente o linealmente dependiente.

\item Determine si el conjunto
\[
\{(1,0,1),(0,1,2),(2,1,4)\}
\]
es linealmente independiente o linealmente dependiente.

\item Determine los valores de \(\lambda\in\mathbb{R}\) para los cuales
\[
\{(1,0,\lambda),(2,1,1),(0,\lambda,1)\}
\]
es una base de \(\mathbb{R}^3\).

\item Determine una base y la dimensión del subespacio
\[
W=\{(x,y,z,t)\in\mathbb{R}^4:x+y-z+t=0,\; 2x-y+z=0\}.
\]

\item Sea
\[
V=\left\langle
(1,0,1,0),(2,1,0,1),(3,1,1,1)
\right\rangle.
\]
Determine una base para \(V\) y calcule \(\dim(V)\).

\item Determine si \(T:\mathbb{R}^2\to\mathbb{R}^3\), definida por
\[
T(x,y)=(x+y,2x-y,x),
\]
es una transformación lineal.

\item Sea \(T:\mathbb{R}^2\to\mathbb{R}^3\), dada por
\[
T(x,y)=(x-y,2x+y,3y).
\]
Determine:
\begin{enumerate}[label=\alph*)]
    \item \(T(1,2)\).
    \item \(T(2,-1)\).
    \item \(T(a,b)\).
\end{enumerate}

\item Sea \(T:\mathbb{R}^3\to\mathbb{R}^2\), definida por
\[
T(x,y,z)=(x+y-z,2x-y+z).
\]
Determine \(\ker(T)\), una base para \(\ker(T)\) y su dimensión.

\item Sea \(T:\mathbb{R}^3\to\mathbb{R}^3\), definida por
\[
T(x,y,z)=(x+y,y+z,x+z).
\]
Determine:
\begin{enumerate}[label=\alph*)]
    \item \(\ker(T)\).
    \item \(\operatorname{Im}(T)\).
    \item \(\dim(\operatorname{Im}(T))\).
\end{enumerate}

\item Sea \(T:\mathbb{R}^3\to\mathbb{R}^3\), definida por
\[
T(x,y,z)=(x+2y-z,2x+y+z,x-y+2z).
\]
Determine si \(T\) es inyectiva, sobreyectiva, biyectiva y si corresponde a un isomorfismo.
Justifique usando el Teorema Núcleo--Imagen.

\end{enumerate}

\end{ejercicios}

\begin{soluciones}

\begin{solbox}{Solución ejercicio 1}
Se quiere verificar si:
\[
W=\{(x,y,z)\in\mathbb{R}^3:x-2y+z=0\}
\]
es un subespacio de \(\mathbb{R}^3\).

Para esto usamos el criterio de subespacio: debe contener al vector cero, ser cerrado para la suma y ser cerrado para la ponderación escalar.

Primero, verificamos si el vector cero pertenece a \(W\):
\[
0-2\cdot 0+0=0.
\]
Por lo tanto:
\[
(0,0,0)\in W.
\]

Ahora tomemos dos vectores cualesquiera de \(W\):
\[
u=(x_1,y_1,z_1),\qquad v=(x_2,y_2,z_2).
\]
Como ambos pertenecen a \(W\), cumplen:
\[
x_1-2y_1+z_1=0,
\]
\[
x_2-2y_2+z_2=0.
\]

Estudiamos la suma:
\[
u+v=(x_1+x_2,y_1+y_2,z_1+z_2).
\]
Evaluamos la condición del conjunto:
\[
(x_1+x_2)-2(y_1+y_2)+(z_1+z_2).
\]
Ordenando:
\[
(x_1-2y_1+z_1)+(x_2-2y_2+z_2)=0+0=0.
\]
Entonces:
\[
u+v\in W.
\]

Finalmente, sea \(\alpha\in\mathbb{R}\). Si \(u\in W\), entonces:
\[
\alpha u=(\alpha x_1,\alpha y_1,\alpha z_1).
\]
Evaluamos:
\[
\alpha x_1-2\alpha y_1+\alpha z_1
=\alpha(x_1-2y_1+z_1)=\alpha\cdot 0=0.
\]
Por lo tanto:
\[
\alpha u\in W.
\]

\begin{conclusionbox}
En conclusión, \(W\) es un subespacio de \(\mathbb{R}^3\).
\end{conclusionbox}
\end{solbox}

\begin{solbox}{Solución ejercicio 2}
Se quiere determinar si:
\[
W=\{(x,y,z)\in\mathbb{R}^3:x+y=1\}
\]
es un subespacio de \(\mathbb{R}^3\).

Para que un conjunto sea subespacio, una condición necesaria es que contenga al vector cero:
\[
(0,0,0).
\]
Evaluamos la condición \(x+y=1\) en el vector cero:
\[
0+0=0.
\]
Pero:
\[
0\neq 1.
\]
Por lo tanto:
\[
(0,0,0)\notin W.
\]

Como falla esta condición básica, no es necesario seguir revisando cerradura por suma o por ponderación escalar.

\begin{conclusionbox}
En conclusión, \(W\) no es un subespacio de \(\mathbb{R}^3\), ya que no contiene al vector cero.
\end{conclusionbox}
\end{solbox}

\begin{solbox}{Solución ejercicio 3}
Se quiere verificar si:
\[
W=\left\{
\begin{pmatrix}
a&b\\
c&d
\end{pmatrix}\in M_{\mathbb{R}}(2):a+d=0
\right\}
\]
es un subespacio de \(M_{\mathbb{R}}(2)\).

La condición \(a+d=0\) indica que la suma de los elementos de la diagonal principal debe ser cero.

Primero verificamos si la matriz cero pertenece a \(W\):
\[
\begin{pmatrix}
0&0\\
0&0
\end{pmatrix}.
\]
En este caso:
\[
0+0=0.
\]
Por lo tanto, la matriz cero pertenece a \(W\).

Ahora tomemos:
\[
A=
\begin{pmatrix}
a&b\\
c&d
\end{pmatrix}\in W,
\qquad
B=
\begin{pmatrix}
e&f\\
g&h
\end{pmatrix}\in W.
\]
Entonces:
\[
a+d=0,
\qquad
e+h=0.
\]

Estudiamos la suma:
\[
A+B=
\begin{pmatrix}
a+e&b+f\\
c+g&d+h
\end{pmatrix}.
\]
La suma de su diagonal principal es:
\[
(a+e)+(d+h).
\]
Reordenando:
\[
(a+d)+(e+h)=0+0=0.
\]
Luego:
\[
A+B\in W.
\]

Finalmente, sea \(\alpha\in\mathbb{R}\). Entonces:
\[
\alpha A=
\begin{pmatrix}
\alpha a&\alpha b\\
\alpha c&\alpha d
\end{pmatrix}.
\]
La suma de su diagonal principal es:
\[
\alpha a+\alpha d=\alpha(a+d)=\alpha\cdot 0=0.
\]
Así:
\[
\alpha A\in W.
\]

\begin{conclusionbox}
En conclusión, \(W\) es un subespacio de \(M_{\mathbb{R}}(2)\).
\end{conclusionbox}
\end{solbox}

\begin{solbox}{Solución ejercicio 4}
Se quiere verificar si:
\[
W=\{a_0+a_1x+a_2x^2\in\mathbb{R}_2[x]:a_0-a_1+a_2=0\}
\]
es un subespacio de \(\mathbb{R}_2[x]\).

Primero verificamos si el polinomio cero pertenece a \(W\):
\[
0+0x+0x^2.
\]
Sus coeficientes son:
\[
a_0=0,\qquad a_1=0,\qquad a_2=0.
\]
Entonces:
\[
a_0-a_1+a_2=0-0+0=0.
\]
Por lo tanto, el polinomio cero pertenece a \(W\).

Ahora tomemos:
\[
p(x)=a_0+a_1x+a_2x^2\in W,
\]
\[
q(x)=b_0+b_1x+b_2x^2\in W.
\]
Entonces:
\[
a_0-a_1+a_2=0,
\]
\[
b_0-b_1+b_2=0.
\]

Estudiamos la suma:
\[
p(x)+q(x)=(a_0+b_0)+(a_1+b_1)x+(a_2+b_2)x^2.
\]
Evaluamos la condición:
\[
(a_0+b_0)-(a_1+b_1)+(a_2+b_2).
\]
Agrupando:
\[
(a_0-a_1+a_2)+(b_0-b_1+b_2)=0+0=0.
\]
Luego:
\[
p(x)+q(x)\in W.
\]

Sea \(\alpha\in\mathbb{R}\). Entonces:
\[
\alpha p(x)=\alpha a_0+\alpha a_1x+\alpha a_2x^2.
\]
Evaluamos:
\[
\alpha a_0-\alpha a_1+\alpha a_2
=\alpha(a_0-a_1+a_2)=0.
\]
Por lo tanto:
\[
\alpha p(x)\in W.
\]

\begin{conclusionbox}
En conclusión, \(W\) es un subespacio de \(\mathbb{R}_2[x]\).
\end{conclusionbox}
\end{solbox}

\begin{solbox}{Solución ejercicio 5}
Se tiene:
\[
W=\{(x,y,z)\in\mathbb{R}^3:x+y-z=0\}.
\]
Para encontrar un sistema generador, debemos escribir cualquier vector de \(W\) como combinación lineal de algunos vectores.

Partimos desde la ecuación:
\[
x+y-z=0.
\]
Despejamos una variable en función de las otras. En este caso:
\[
x=-y+z.
\]
Como \(y\) y \(z\) quedan libres, tomamos:
\[
y=s,\qquad z=t.
\]
Entonces:
\[
x=-s+t.
\]
Por lo tanto:
\[
(x,y,z)=(-s+t,s,t).
\]

Ahora separamos según cada parámetro:
\[
(-s+t,s,t)=(-s,s,0)+(t,0,t).
\]
Factorizando:
\[
(-s,s,0)=s(-1,1,0),
\]
\[
(t,0,t)=t(1,0,1).
\]
Así:
\[
(x,y,z)=s(-1,1,0)+t(1,0,1).
\]
Por lo tanto:
\[
W=\langle\{(-1,1,0),(1,0,1)\}\rangle.
\]

\begin{conclusionbox}
En conclusión, un sistema generador para \(W\) es:
\[
\{(-1,1,0),(1,0,1)\}.
\]
\end{conclusionbox}
\end{solbox}

\begin{solbox}{Solución ejercicio 6}
Se tiene:
\[
W=\{(x,y,z,t)\in\mathbb{R}^4:x-y+t=0\}.
\]
La idea es expresar el vector general \((x,y,z,t)\) en función de parámetros libres.

Partimos con:
\[
x-y+t=0.
\]
Despejamos \(x\):
\[
x=y-t.
\]
Las variables \(y,z,t\) quedan libres. Tomamos:
\[
y=s,\qquad z=r,\qquad t=q.
\]
Entonces:
\[
x=s-q.
\]
Por lo tanto:
\[
(x,y,z,t)=(s-q,s,r,q).
\]

Ahora separamos por parámetros:
\[
(s-q,s,r,q)=(s,s,0,0)+(0,0,r,0)+(-q,0,0,q).
\]
Factorizando:
\[
(s,s,0,0)=s(1,1,0,0),
\]
\[
(0,0,r,0)=r(0,0,1,0),
\]
\[
(-q,0,0,q)=q(-1,0,0,1).
\]
Así:
\[
(x,y,z,t)=s(1,1,0,0)+r(0,0,1,0)+q(-1,0,0,1).
\]

\begin{conclusionbox}
En conclusión, un sistema generador para \(W\) es:
\[
\{(1,1,0,0),(0,0,1,0),(-1,0,0,1)\}.
\]
\end{conclusionbox}
\end{solbox}

\begin{solbox}{Solución ejercicio 7}
Se tiene:
\[
W=\left\{
\begin{pmatrix}
a&b\\
c&d
\end{pmatrix}\in M_{\mathbb{R}}(2):a=d
\right\}.
\]
La condición \(a=d\) dice que los elementos de la diagonal principal son iguales.

Entonces toda matriz de \(W\) tiene la forma:
\[
\begin{pmatrix}
a&b\\
c&a
\end{pmatrix}.
\]
Ahora separamos la matriz según sus parámetros \(a,b,c\):
\[
\begin{pmatrix}
a&b\\
c&a
\end{pmatrix}
=
\begin{pmatrix}
a&0\\
0&a
\end{pmatrix}
+
\begin{pmatrix}
0&b\\
0&0
\end{pmatrix}
+
\begin{pmatrix}
0&0\\
c&0
\end{pmatrix}.
\]
Factorizando:
\[
\begin{pmatrix}
a&0\\
0&a
\end{pmatrix}
=
a\begin{pmatrix}
1&0\\
0&1
\end{pmatrix},
\]
\[
\begin{pmatrix}
0&b\\
0&0
\end{pmatrix}
=
b\begin{pmatrix}
0&1\\
0&0
\end{pmatrix},
\]
\[
\begin{pmatrix}
0&0\\
c&0
\end{pmatrix}
=
c\begin{pmatrix}
0&0\\
1&0
\end{pmatrix}.
\]
Por lo tanto:
\[
W=
\left\langle
\left\{
\begin{pmatrix}1&0\\0&1\end{pmatrix},
\begin{pmatrix}0&1\\0&0\end{pmatrix},
\begin{pmatrix}0&0\\1&0\end{pmatrix}
\right\}
\right\rangle.
\]

\begin{conclusionbox}
En conclusión, un sistema generador para \(W\) es:
\[
\left\{
\begin{pmatrix}1&0\\0&1\end{pmatrix},
\begin{pmatrix}0&1\\0&0\end{pmatrix},
\begin{pmatrix}0&0\\1&0\end{pmatrix}
\right\}.
\]
\end{conclusionbox}
\end{solbox}

\begin{solbox}{Solución ejercicio 8}
Sea:
\[
p(x)=a_0+a_1x+a_2x^2.
\]
La condición del conjunto es:
\[
p(1)=0.
\]
Evaluamos:
\[
p(1)=a_0+a_1+a_2.
\]
Entonces:
\[
a_0+a_1+a_2=0.
\]
Despejamos:
\[
a_0=-a_1-a_2.
\]
Las variables \(a_1\) y \(a_2\) quedan libres. Tomamos:
\[
a_1=s,\qquad a_2=t.
\]
Entonces:
\[
a_0=-s-t.
\]
Por lo tanto:
\[
p(x)=(-s-t)+sx+tx^2.
\]
Separamos por parámetros:
\[
p(x)=s(-1+x)+t(-1+x^2).
\]
Así, todo polinomio de \(W\) se puede escribir como combinación lineal de \(-1+x\) y \(-1+x^2\).

\begin{conclusionbox}
En conclusión, un sistema generador para \(W\) es:
\[
\{-1+x,\,-1+x^2\}.
\]
\end{conclusionbox}
\end{solbox}

\begin{solbox}{Solución ejercicio 9}
Queremos determinar si:
\[
\{(1,2),(2,4)\}
\]
es linealmente independiente o dependiente.

Para eso estudiamos:
\[
\alpha(1,2)+\beta(2,4)=(0,0).
\]
Al sumar componente a componente:
\[
(\alpha+2\beta,2\alpha+4\beta)=(0,0).
\]
Por lo tanto:
\[
\alpha+2\beta=0,
\]
\[
2\alpha+4\beta=0.
\]
La segunda ecuación es simplemente dos veces la primera:
\[
2(\alpha+2\beta)=0.
\]
Por lo tanto, no entrega nueva información. Esto permite soluciones no triviales.

Por ejemplo, si tomamos:
\[
\beta=1,
\]
entonces:
\[
\alpha+2=0 \Rightarrow \alpha=-2.
\]
Así:
\[
-2(1,2)+1(2,4)=(0,0).
\]
Como encontramos una combinación no trivial que da el vector cero, el conjunto es linealmente dependiente.

\begin{conclusionbox}
En conclusión, el conjunto es linealmente dependiente.
\end{conclusionbox}
\end{solbox}

\begin{solbox}{Solución ejercicio 10}
Queremos estudiar:
\[
\{(1,0,1),(0,1,2),(2,1,4)\}.
\]
Planteamos la combinación lineal igual al vector cero:
\[
\alpha(1,0,1)+\beta(0,1,2)+\gamma(2,1,4)=(0,0,0).
\]
Sumando componente a componente:
\[
(\alpha+2\gamma,\beta+\gamma,\alpha+2\beta+4\gamma)=(0,0,0).
\]
Entonces:
\[
\alpha+2\gamma=0,
\]
\[
\beta+\gamma=0,
\]
\[
\alpha+2\beta+4\gamma=0.
\]
De la primera ecuación:
\[
\alpha=-2\gamma.
\]
De la segunda:
\[
\beta=-\gamma.
\]
Reemplazamos en la tercera:
\[
-2\gamma+2(-\gamma)+4\gamma=0.
\]
Esto se cumple para cualquier valor de \(\gamma\). Por lo tanto, podemos elegir un valor no nulo. Por ejemplo, si \(\gamma=1\), entonces:
\[
\alpha=-2,\qquad \beta=-1.
\]
Con esto tenemos una combinación no trivial que da el vector cero.

\begin{conclusionbox}
En conclusión, el conjunto es linealmente dependiente.
\end{conclusionbox}
\end{solbox}

\begin{solbox}{Solución ejercicio 11}
Para que el conjunto:
\[
\{(1,0,\lambda),(2,1,1),(0,\lambda,1)\}
\]
sea base de \(\mathbb{R}^3\), debe tener tres vectores linealmente independientes. Como son tres vectores en \(\mathbb{R}^3\), podemos formar una matriz con ellos y exigir que su determinante sea distinto de cero.

Tomamos los vectores como columnas:
\[
A=
\begin{pmatrix}
1&2&0\\
0&1&\lambda\\
\lambda&1&1
\end{pmatrix}.
\]
El conjunto será base si:
\[
|A|\neq 0.
\]

Calculamos el determinante expandiendo por la primera fila:
\[
|A|
=
1\begin{vmatrix}
1&\lambda\\
1&1
\end{vmatrix}
-
2\begin{vmatrix}
0&\lambda\\
\lambda&1
\end{vmatrix}
+
0\begin{vmatrix}
0&1\\
\lambda&1
\end{vmatrix}.
\]
Entonces:
\[
|A|=1(1-\lambda)-2(0-\lambda^2)+0.
\]
Por lo tanto:
\[
|A|=1-\lambda+2\lambda^2.
\]
Ordenando:
\[
|A|=2\lambda^2-\lambda+1.
\]

Ahora estudiamos cuándo este determinante podría ser cero:
\[
2\lambda^2-\lambda+1=0.
\]
El discriminante es:
\[
\Delta=(-1)^2-4(2)(1)=1-8=-7.
\]
Como:
\[
\Delta<0,
\]
la ecuación no tiene soluciones reales. Además, como el coeficiente de \(\lambda^2\) es positivo, la expresión nunca se anula en \(\mathbb{R}\).

Entonces:
\[
2\lambda^2-\lambda+1\neq 0,\quad \forall \lambda\in\mathbb{R}.
\]

\begin{conclusionbox}
En conclusión, el conjunto es base de \(\mathbb{R}^3\) para todo \(\lambda\in\mathbb{R}\).
\end{conclusionbox}
\end{solbox}

\begin{solbox}{Solución ejercicio 12}
El subespacio está dado por las condiciones:
\[
x+y-z+t=0,\qquad 2x-y+z=0.
\]
La idea es encontrar todos los vectores \((x,y,z,t)\) que cumplen ambas ecuaciones. Para ello resolvemos el sistema homogéneo asociado:
\[
\begin{pmatrix}
1&1&-1&1&|&0\\
2&-1&1&0&|&0
\end{pmatrix}.
\]

Eliminamos la \(x\) de la segunda fila:
\[
f_2\leftarrow f_2-2f_1.
\]
Entonces:
\[
f_2=(2,-1,1,0)-2(1,1,-1,1)=(0,-3,3,-2).
\]
Por lo tanto:
\[
\begin{pmatrix}
1&1&-1&1&|&0\\
0&-3&3&-2&|&0
\end{pmatrix}.
\]

La segunda fila representa:
\[
-3y+3z-2t=0.
\]
Despejamos \(y\):
\[
-3y=-3z+2t
\]
\[
y=z-\frac{2}{3}t.
\]

Ahora usamos la primera ecuación:
\[
x+y-z+t=0.
\]
Despejamos \(x\):
\[
x=-y+z-t.
\]
Reemplazando \(y=z-\frac{2}{3}t\):
\[
x=-\left(z-\frac{2}{3}t\right)+z-t.
\]
Entonces:
\[
x=-z+\frac{2}{3}t+z-t=-\frac{1}{3}t.
\]

Así:
\[
x=-\frac{1}{3}t,\qquad y=z-\frac{2}{3}t.
\]
Como \(z\) y \(t\) quedan libres, los usamos como parámetros. Para evitar fracciones, tomamos:
\[
z=s,\qquad t=3r.
\]
Luego:
\[
x=-r,\qquad y=s-2r,\qquad z=s,\qquad t=3r.
\]
Por lo tanto:
\[
(x,y,z,t)=(-r,s-2r,s,3r).
\]

Ahora separamos el vector según cada parámetro:
\[
(x,y,z,t)=(0,s,s,0)+(-r,-2r,0,3r).
\]
Factorizando:
\[
(x,y,z,t)=s(0,1,1,0)+r(-1,-2,0,3).
\]

Esto significa que todo vector de \(W\) se puede escribir como combinación lineal de:
\[
(0,1,1,0)\quad \text{y}\quad (-1,-2,0,3).
\]
Además, estos dos vectores no son múltiplos entre sí, por lo que son linealmente independientes.

\begin{conclusionbox}
En conclusión, una base para \(W\) es:
\[
\{(0,1,1,0),(-1,-2,0,3)\},
\]
y:
\[
\dim(W)=2.
\]
\end{conclusionbox}
\end{solbox}

\begin{solbox}{Solución ejercicio 13}
El espacio \(V\) está generado por:
\[
v_1=(1,0,1,0),\qquad
v_2=(2,1,0,1),\qquad
v_3=(3,1,1,1).
\]
Para encontrar una base, buscamos eliminar vectores redundantes, es decir, vectores que se puedan escribir como combinación lineal de los otros.

Estudiamos si los tres generadores son linealmente independientes:
\[
\alpha v_1+\beta v_2+\gamma v_3=(0,0,0,0).
\]
Esto es:
\[
\alpha(1,0,1,0)+\beta(2,1,0,1)+\gamma(3,1,1,1)=(0,0,0,0).
\]
Al sumar componente a componente:
\[
(\alpha+2\beta+3\gamma,\beta+\gamma,\alpha+\gamma,\beta+\gamma)=(0,0,0,0).
\]
Por lo tanto:
\[
\alpha+2\beta+3\gamma=0,
\]
\[
\beta+\gamma=0,
\]
\[
\alpha+\gamma=0,
\]
\[
\beta+\gamma=0.
\]
De \(\beta+\gamma=0\), se obtiene:
\[
\beta=-\gamma.
\]
De \(\alpha+\gamma=0\), se obtiene:
\[
\alpha=-\gamma.
\]
Reemplazando en la primera ecuación:
\[
-\gamma+2(-\gamma)+3\gamma=0.
\]
Esto se cumple para todo \(\gamma\), por lo que existen soluciones no triviales. Entonces los tres vectores son linealmente dependientes.

Ahora observamos directamente que:
\[
v_1+v_2=(1,0,1,0)+(2,1,0,1)=(3,1,1,1)=v_3.
\]
Por lo tanto, \(v_3\) no aporta una nueva dirección al espacio generado.

Nos quedamos con:
\[
v_1=(1,0,1,0),\qquad v_2=(2,1,0,1).
\]
Estos dos vectores no son múltiplos entre sí, por lo que son linealmente independientes. Además, generan el mismo espacio \(V\), porque \(v_3=v_1+v_2\).

\begin{conclusionbox}
En conclusión, una base para \(V\) es:
\[
\{(1,0,1,0),(2,1,0,1)\},
\]
y:
\[
\dim(V)=2.
\]
\end{conclusionbox}
\end{solbox}

\begin{solbox}{Solución ejercicio 14}
Se quiere verificar si:
\[
T:\mathbb{R}^2\to\mathbb{R}^3,\qquad T(x,y)=(x+y,2x-y,x)
\]
es transformación lineal.

Para ello debemos probar dos propiedades:

\[
T(u+v)=T(u)+T(v)
\]
y
\[
T(\alpha u)=\alpha T(u).
\]

Sea:
\[
u=(x_1,y_1),\qquad v=(x_2,y_2).
\]
Entonces:
\[
u+v=(x_1+x_2,y_1+y_2).
\]
Aplicamos \(T\):
\[
T(u+v)=T(x_1+x_2,y_1+y_2).
\]
Por definición de \(T\):
\[
T(u+v)=((x_1+x_2)+(y_1+y_2),2(x_1+x_2)-(y_1+y_2),x_1+x_2).
\]
Ordenando:
\[
T(u+v)=(x_1+y_1+x_2+y_2,2x_1-y_1+2x_2-y_2,x_1+x_2).
\]
Esto se puede separar como:
\[
(x_1+y_1,2x_1-y_1,x_1)+(x_2+y_2,2x_2-y_2,x_2).
\]
Por lo tanto:
\[
T(u+v)=T(u)+T(v).
\]

Ahora sea \(\alpha\in\mathbb{R}\). Entonces:
\[
\alpha u=(\alpha x_1,\alpha y_1).
\]
Aplicamos \(T\):
\[
T(\alpha u)=T(\alpha x_1,\alpha y_1).
\]
Por definición:
\[
T(\alpha u)=(\alpha x_1+\alpha y_1,2\alpha x_1-\alpha y_1,\alpha x_1).
\]
Factorizamos \(\alpha\):
\[
T(\alpha u)=\alpha(x_1+y_1,2x_1-y_1,x_1).
\]
Por lo tanto:
\[
T(\alpha u)=\alpha T(u).
\]

\begin{conclusionbox}
En conclusión, \(T\) es una transformación lineal.
\end{conclusionbox}
\end{solbox}

\begin{solbox}{Solución ejercicio 15}
La transformación está dada por:
\[
T(x,y)=(x-y,2x+y,3y).
\]
Para evaluar una transformación, reemplazamos las coordenadas del vector en la regla de \(T\).

\textbf{a)} Para \((1,2)\):
\[
T(1,2)=(1-2,2(1)+2,3(2)).
\]
Calculando componente a componente:
\[
T(1,2)=(-1,4,6).
\]

\textbf{b)} Para \((2,-1)\):
\[
T(2,-1)=(2-(-1),2(2)+(-1),3(-1)).
\]
Entonces:
\[
T(2,-1)=(3,3,-3).
\]

\textbf{c)} Para \((a,b)\):
\[
T(a,b)=(a-b,2a+b,3b).
\]
Aquí no se reemplazan valores numéricos, porque se pide la imagen de un vector general.

\begin{conclusionbox}
En conclusión:
\[
T(1,2)=(-1,4,6),\qquad T(2,-1)=(3,3,-3),
\]
y:
\[
T(a,b)=(a-b,2a+b,3b).
\]
\end{conclusionbox}
\end{solbox}

\begin{solbox}{Solución ejercicio 16}
Por definición, el núcleo está formado por todos los vectores del dominio que la transformación envía al vector cero:
\[
\ker(T)=\{(x,y,z)\in\mathbb{R}^3:T(x,y,z)=(0,0)\}.
\]
Como:
\[
T(x,y,z)=(x+y-z,2x-y+z),
\]
debemos resolver:
\[
(x+y-z,2x-y+z)=(0,0).
\]
Esto genera el sistema:
\[
x+y-z=0,
\]
\[
2x-y+z=0.
\]

Sumamos ambas ecuaciones para eliminar \(y\) y \(z\):
\[
(x+y-z)+(2x-y+z)=0+0.
\]
Así:
\[
3x=0.
\]
Por lo tanto:
\[
x=0.
\]
Reemplazamos \(x=0\) en la primera ecuación:
\[
0+y-z=0.
\]
Entonces:
\[
y=z.
\]
Como \(z\) queda libre, tomamos:
\[
z=t.
\]
Luego:
\[
y=t,\qquad x=0.
\]
Por lo tanto:
\[
(x,y,z)=(0,t,t).
\]
Factorizando:
\[
(x,y,z)=t(0,1,1).
\]
Así:
\[
\ker(T)=\langle\{(0,1,1)\}\rangle.
\]

\begin{conclusionbox}
En conclusión:
\[
\ker(T)=\{(0,t,t):t\in\mathbb{R}\},
\]
una base para \(\ker(T)\) es:
\[
\{(0,1,1)\},
\]
y:
\[
\dim(\ker T)=1.
\]
\end{conclusionbox}
\end{solbox}

\begin{solbox}{Solución ejercicio 17}
La transformación está dada por:
\[
T(x,y,z)=(x+y,y+z,x+z).
\]

\textbf{a)} Para calcular el núcleo, resolvemos:
\[
T(x,y,z)=(0,0,0).
\]
Entonces:
\[
(x+y,y+z,x+z)=(0,0,0).
\]
Esto genera el sistema:
\[
x+y=0,
\]
\[
y+z=0,
\]
\[
x+z=0.
\]
De la primera ecuación:
\[
x=-y.
\]
De la segunda:
\[
z=-y.
\]
Reemplazamos en la tercera:
\[
x+z=0.
\]
Como \(x=-y\) y \(z=-y\):
\[
-y-y=0.
\]
Entonces:
\[
-2y=0 \Rightarrow y=0.
\]
Luego:
\[
x=0,\qquad z=0.
\]
Por lo tanto:
\[
\ker(T)=\{(0,0,0)\}.
\]

\textbf{b)} Para encontrar la imagen, calculamos la imagen de la base canónica de \(\mathbb{R}^3\):
\[
e_1=(1,0,0),\qquad e_2=(0,1,0),\qquad e_3=(0,0,1).
\]
Entonces:
\[
T(1,0,0)=(1,0,1),
\]
\[
T(0,1,0)=(1,1,0),
\]
\[
T(0,0,1)=(0,1,1).
\]
Por lo tanto:
\[
\operatorname{Im}(T)=
\left\langle
\{(1,0,1),(1,1,0),(0,1,1)\}
\right\rangle.
\]

\textbf{c)} Como:
\[
\ker(T)=\{(0,0,0)\},
\]
se tiene:
\[
\dim(\ker T)=0.
\]
Aplicamos el Teorema Núcleo--Imagen:
\[
\dim(\ker T)+\dim(\operatorname{Im}T)=\dim(\mathbb{R}^3).
\]
Reemplazando:
\[
0+\dim(\operatorname{Im}T)=3.
\]
Por lo tanto:
\[
\dim(\operatorname{Im}T)=3.
\]

\begin{conclusionbox}
En conclusión:
\[
\ker(T)=\{(0,0,0)\},
\]
\[
\operatorname{Im}(T)=
\left\langle
\{(1,0,1),(1,1,0),(0,1,1)\}
\right\rangle,
\]
y:
\[
\dim(\operatorname{Im}T)=3.
\]
\end{conclusionbox}
\end{solbox}

\begin{solbox}{Solución ejercicio 18}
Para estudiar inyectividad y sobreyectividad, comenzamos calculando el núcleo. Si el núcleo contiene solo al vector cero, entonces \(T\) es inyectiva.

Resolvemos:
\[
T(x,y,z)=(0,0,0).
\]
Como:
\[
T(x,y,z)=(x+2y-z,2x+y+z,x-y+2z),
\]
obtenemos el sistema:
\[
x+2y-z=0,
\]
\[
2x+y+z=0,
\]
\[
x-y+2z=0.
\]

Escribimos la matriz aumentada asociada:
\[
\begin{pmatrix}
1&2&-1&|&0\\
2&1&1&|&0\\
1&-1&2&|&0
\end{pmatrix}.
\]
Eliminamos la \(x\) de la segunda fila:
\[
f_2\leftarrow f_2-2f_1.
\]
Entonces:
\[
f_2=(2,1,1)-2(1,2,-1)=(0,-3,3).
\]
Eliminamos la \(x\) de la tercera fila:
\[
f_3\leftarrow f_3-f_1.
\]
Entonces:
\[
f_3=(1,-1,2)-(1,2,-1)=(0,-3,3).
\]
Por lo tanto:
\[
\begin{pmatrix}
1&2&-1&|&0\\
0&-3&3&|&0\\
0&-3&3&|&0
\end{pmatrix}
\sim
\begin{pmatrix}
1&2&-1&|&0\\
0&-3&3&|&0\\
0&0&0&|&0
\end{pmatrix}.
\]

De la segunda fila:
\[
-3y+3z=0.
\]
Por lo tanto:
\[
y=z.
\]
De la primera fila:
\[
x+2y-z=0.
\]
Como \(y=z\), reemplazamos:
\[
x+2z-z=0.
\]
Entonces:
\[
x+z=0 \Rightarrow x=-z.
\]
Sea \(z=t\). Entonces:
\[
x=-t,\qquad y=t,\qquad z=t.
\]
Así:
\[
(x,y,z)=(-t,t,t)=t(-1,1,1).
\]
Por lo tanto:
\[
\ker(T)=\langle\{(-1,1,1)\}\rangle.
\]
Luego:
\[
\dim(\ker T)=1.
\]

Como el núcleo no es trivial, \(T\) no es inyectiva.

Ahora usamos el Teorema Núcleo--Imagen:
\[
\dim(\ker T)+\dim(\operatorname{Im}T)=\dim(\mathbb{R}^3).
\]
Reemplazando:
\[
1+\dim(\operatorname{Im}T)=3.
\]
Entonces:
\[
\dim(\operatorname{Im}T)=2.
\]
Como la imagen tiene dimensión \(2\) y el codominio \(\mathbb{R}^3\) tiene dimensión \(3\), la imagen no puede ser todo \(\mathbb{R}^3\). Por lo tanto, \(T\) no es sobreyectiva.

\begin{conclusionbox}
En conclusión, \(T\) no es inyectiva, no es sobreyectiva y, por lo tanto, no es biyectiva.
Como \(T\) no es biyectiva, entonces \(T\) no corresponde a un isomorfismo.
\end{conclusionbox}
\end{solbox}

\end{soluciones}

\end{document}
